\documentclass[aps,prd,twocolumn,nofootinbib,superscriptaddress]{revtex4-2}

\usepackage[T1]{fontenc}
\usepackage[utf8]{inputenc}
\usepackage{graphicx}
\usepackage{amsmath,amssymb,amsthm}
\usepackage{bm}

\usepackage{orcidlink}
\usepackage{hyperref}

% Define theorem environments for the formal section
\newtheorem{theorem}{Theorem}
\newtheorem{definition}{Definition}
\newtheorem{lemma}{Lemma}

\begin{document}

\title{Causal Rigidity and the Single-Unit Universe: \\ Integrating the Alexandrov-Zeeman and Unruh clock scales}

\author{Karl Svozil\,\orcidlink{0000-0001-6554-2802}}
\email{karl.svozil@tuwien.ac.at}
\homepage{http://tph.tuwien.ac.at/~svozil}

\affiliation{Institute for Theoretical Physics,
TU Wien,
Wiedner Hauptstrasse 8-10/136,
1040 Vienna,  Austria}

\date{\today}

\begin{abstract}
We unify two complementary viewpoints on relativistic spacetime and the counting of fundamental constants. Operationally, Matsas, Pleitez, Saa, and Vanzella (MPSV) have recently argued that relativistic spacetime requires only a single fundamental dimensional constant, demonstrating this via a clock-based protocol that measures spatial distance without light rays. Mathematically, theorems due to Alexandrov and Zeeman demonstrate that the light-cone structure determines the spacetime geometry only up to a conformal factor. We show that these approaches are mutually reinforcing: the Alexandrov-Zeeman theorems establish the rigid conformal structure of spacetime, while the ``bona fide clock'' required by MPSV serves the necessary mathematical role of breaking the dilation symmetry. We provide a formal derivation proving that the normalization of a single clock worldline is sufficient to select a unique metric from the conformal class, thereby clarifying that the number of fundamental constants is exactly one.
\end{abstract}

\maketitle

\section{Introduction}

The question ``How many fundamental dimensional constants does physics require?'' has been a subject of enduring debate, most famously in the ``trialogue'' between Duff, Okun, and Veneziano (DOV)~\cite{DuffOkunVeneziano}. While there is consensus that only dimensionless quantities are ultimately observable, disagreements persist regarding the number of independent \emph{standards} (or base units) necessary to describe dimensionful observables.

Recently, Matsas, Pleitez, Saa, and Vanzella (MPSV)~\cite{matsas-2024} revisited this issue from a \emph{spacetime-based} perspective. They argue that once a relativistic spacetime is assumed as the starting point, the units required to construct that spacetime are sufficient to express all observables. Specifically, they conclude that in a relativistic spacetime, the number of fundamental dimensional constants is \emph{one}.

Parallel to this metrological debate, the mathematical relativity community has long established---via theorems by Alexandrov~\cite{alex3} and Zeeman~\cite{Zeeman}---that the null (light-cone) structure of spacetime is incredibly rigid. These theorems state that any map preserving the causal order is a Lorentz transformation up to a global dilation (scale factor).

In this paper, we unify these two perspectives. In Sec.~\ref{sec:Operational}, we review the operational argument provided by MPSV. In Sec.~\ref{sec:Geometric}, we review the geometric constraints imposed by causal automorphisms. Finally, in Sec.~\ref{sec:Formal}, we provide a formal derivation demonstrating that the ``bona fide clock'' of MPSV is the precise mathematical object required to break the conformal symmetry left open by incidence geometry.

\section{The Operational Perspective: Measuring Space with Time}
\label{sec:Operational}

MPSV argue that to construct a relativistic spacetime physically, one requires specific apparatuses. Unlike Galilean spacetime, which requires both rulers and clocks (because space and time are decoupled), relativistic spacetime requires only ``bona fide clocks.''

To prove that spatial rulers are redundant, MPSV utilize a protocol attributed to William George Unruh. Consider three inertial clocks $C_1, C_2, C_3$:
\begin{enumerate}
    \item $C_1$ travels from an observer $A$ to $B$.
    \item $C_2$ travels from $B$ back to $A$.
    \item $C_3$ remains at rest with $A$.
\end{enumerate}
If $\tau_1$ and $\tau_2$ are the proper time durations measured by the traveling clocks, and $\tau_3$ is the duration measured by the stationary clock, the spatial distance $D$ between $A$ and $B$ is given by Heron's formula:
\begin{equation}
D =  \frac{\sqrt{(\tau_3^2 - \tau_1^2 - \tau_2^2)^2 - 4\tau_1^2\tau_2^2}}{2\tau_3}.
\label{eq:Unruh}
\end{equation}
This formula yields a spatial length in units of time (seconds). It demonstrates physically that the spatial metric can be derived entirely from time measurements, rendering the meter a ``disposable'' unit and $c$ a mere conversion factor.

\section{The Geometric Perspective: Alexandrov-Zeeman Theorems}
\label{sec:Geometric}

We now turn to the mathematical structure of Minkowski spacetime. We consider the manifold $M \simeq \mathbb{R}^n$ (with $n \geq 3$), equipped with the standard flat metric $\eta_{ab} = \text{diag}(-1, 1, \dots, 1)$. The defining feature of relativistic geometry is the causal structure, encoded either by the incidence of light cones or the partial causal ordering of events.

A fundamental question in mathematical relativity is the extent to which this causal structure alone determines the geometry. This is resolved by a class of rigidity results known as Alexandrov-Zeeman type theorems.

\subsection{Alexandrov: The Rigidity of Light Cones}
In the 1950s, A.~D.~Alexandrov proved that the geometry of spacetime is encoded in the ``skeleton'' of light rays~\cite{alex3}. Let $C_x$ denote the light cone (null cone) at an event $x \in M$.

\begin{theorem}[Alexandrov]
Let $f: M \to M$ be a bijection of the spacetime manifold ($n \ge 3$). If $f$ maps null lines to null lines (i.e., preserves the light-cone structure), then $f$ is a linear conformal automorphism.
\end{theorem}
This result implies that mere incidence---knowing which events can be connected by a light ray---fixes the linear structure of spacetime~\cite{lester}.

\subsection{Zeeman: Causality Implies the Lorentz Group}
In 1964, E.~C.~Zeeman strengthened this perspective by focusing on the partial ordering of events~\cite{Zeeman}. We say $x \prec y$ if a signal can travel from $x$ to $y$ (with speed $v \leq c$).

\begin{theorem}[Zeeman]
Let $f: M \to M$ be a bijection that preserves the causal order (i.e., $x \prec y \iff f(x) \prec f(y)$). Then $f$ is generated by:
\begin{enumerate}
    \item An orthochronous Lorentz transformation $\Lambda \in O^\uparrow(1, n-1)$,
    \item A spacetime translation $a$, and
    \item A global dilation (scaling) $x \mapsto \lambda x$ with $\lambda > 0$.
\end{enumerate}
\end{theorem}

\subsection{The Conformal Ambiguity}
Both theorems lead to the same physical conclusion: the causal structure determines the metric $g_{ab}$ up to a \emph{conformal factor} $\Omega(x)$. Under the strict conditions of mapping the whole of Minkowski space to itself, this factor becomes a global constant, $\Omega(x) = \lambda$.
\begin{equation}
    g'_{ab} = \lambda^2 \eta_{ab}.
\end{equation}
This leads to a profound physical consequence: \textbf{A universe defined solely by light rays (or causal order) cannot distinguish a second from a millennium.} The transformation $t \to 1000 t$ is a valid symmetry of the causal structure. The ``shape'' of physics is preserved, but the duration is not. To distinguish time scales, one must break this conformal symmetry.

\section{Formal Unification: The Clock as Symmetry Breaker}
\label{sec:Formal}

We now formalize the connection between the MPSV clock and the Alexandrov-Zeeman conformal factor. We show that specifying the normalization of a single clock worldline is sufficient to fix the global scale $\lambda$.

\subsection{Definitions}

\begin{definition}[Conformal Class]
Let $\eta$ be the standard Minkowski metric. The \emph{conformal class} $[\eta]$ determined by the causal structure is the set of metrics $\{\lambda^2 \eta : \lambda > 0\}$.
\end{definition}

\begin{definition}[Bona Fide Clock]
Let $\gamma:\mathbb{R}\to M$ be a smooth, future-directed timelike worldline of the clock. A \emph{bona fide clock} is the pair $(\gamma,\tau)$ where $\tau$ is a parameter on the domain $\mathbb{R}$ such that, with respect to the true physical metric $g_{\star}$, the parameter equals the proper time. That is, for all $s\in\mathbb{R}$:
\begin{equation}
g_{\star}\big(\dot\gamma(s),\dot\gamma(s)\big) = -1,
\label{eq:normalization}
\end{equation}
where dot denotes $d/ds$.
\end{definition}

Physically, this condition enforces that the parameter $\tau$ is synchronized with the spacetime geometry. It asserts that the tangent vector $\dot{\gamma}$, which represents the clock's velocity through spacetime, has a constant, unit magnitude. In other words, for every ``tick'' of the parameter $s$, the system traverses exactly one unit of invariant interval defined by $g_{\star}$. This is what elevates $s$ from an arbitrary coordinate label to a physical measure of duration, thereby fixing the scale of the metric.

\subsection{Uniqueness Theorem}

\begin{theorem}[Single Clock Fixes Scale]
Let $M\cong\mathbb{R}^4$ be equipped with a causal order determining the conformal class $[\eta]$. Let $(\gamma,\tau)$ be a single bona fide clock following an inertial trajectory. Then, the requirement that $\tau$ measures proper time with respect to the physical metric $g_{\star} \in [\eta]$ uniquely determines $g_{\star}$.
\end{theorem}

\begin{proof}
Since the causal structure fixes the metric up to a global scale (via the Alexandrov-Zeeman theorems for $n \ge 3$), the physical metric must take the form:
\begin{equation}
g_{\star} = \lambda^2 \eta, \quad \lambda > 0,
\end{equation}
where $\eta$ is an arbitrary reference metric (e.g., standard Minkowski coordinates).

The bona fide clock condition (Eq.~\ref{eq:normalization}) requires that along the worldline $\gamma$:
\begin{equation}
g_{\star}(\dot{\gamma}, \dot{\gamma}) = -1 \text{ which implies } \lambda^2 \eta(\dot{\gamma}, \dot{\gamma}) = -1.
\end{equation}
Let $-\alpha^2 \equiv \eta(\dot{\gamma}, \dot{\gamma})$ be the squared norm of the clock's tangent vector as measured in the reference metric $\eta$. Since $\gamma$ is timelike, $\alpha^2 > 0$. This quantity $\alpha$ physically encodes the ``tick rate'' of the clock relative to the arbitrary reference coordinates. It represents the real, positive magnitude of the clock's velocity vector through the reference background $\eta$, quantifying the mismatch between the coordinate measuring system and the physical clock.

The equation becomes:
\begin{equation}
\lambda^2 (-\alpha^2) = -1 \implies \lambda = \frac{1}{\alpha}.
\end{equation}
Since $\alpha$ is a determined value based on the coordinate description of the curve in the reference system, $\lambda$ is uniquely determined. Thus, $g_{\star}$ is unique.
\end{proof}

\subsection{The Hierarchy of Metrics}

It is instructive to explicitly distinguish the roles of the metric tensors discussed above. The relationship between the generic conformal scaling and the specific solution $g_{\star}$ can be summarized as follows:

\begin{enumerate}
    \item \textbf{The Reference Metric ($\eta_{ab}$):} An arbitrary mathematical auxiliary (the standard Minkowski metric) that correctly describes the causal incidence (light cones) but possesses no physical scale.
    \item \textbf{The Conformal Class ($g'_{ab}$):} The set of all possible metrics allowed by the Alexandrov-Zeeman theorems. This is a one-parameter family generated by scaling the reference metric:
    \begin{equation}
        g'_{ab} = \lambda^2 \eta_{ab}, \quad \forall \lambda \in (0, \infty).
    \end{equation}
    At this level, the theory cannot distinguish between a second and a millennium; any value of $\lambda$ preserves the causal order.
    \item \textbf{The Physical Metric ($g_{\star}$):} The unique element of the conformal class that is selected by the physical apparatus.
\end{enumerate}

Mathematically, $g_{\star}$ is obtained by evaluating the general form $g'_{ab}$ at the specific scale factor $\lambda_{\text{phys}}$ determined by the clock normalization:
\begin{equation}
    g_{\star} \equiv \left. g'_{ab} \right|_{\lambda = 1/\alpha} = \frac{1}{\alpha^2} \eta_{ab}.
\end{equation}
Thus, while the causal structure generates the candidate list defined by $\eta_{ab}$, the single bona fide clock collapses this list to the unique physical reality $g_{\star}$.

\section{Discussion}

The formal derivation above confirms the operational intuition of MPSV. The ``one fundamental constant'' is the physical realization of the mathematical mechanism required to break the dilation symmetry of incidence geometry.

\begin{itemize}
    \item \textit{Incidence Relations (Alexandrov/Zeeman):} These fix the conformal class $[\eta]$. Physically, they define the ``shape'' of spacetime and the structural role of $c$, but leave the scale undetermined.
    \item \textit{The Clock:} This provides the normalization condition $g(\dot{\gamma}, \dot{\gamma}) = -1$, eliminating the dilation sector.
\end{itemize}

Therefore, the counting of fundamental constants in Special Relativity is exactly one because the symmetry group of the underlying structure is the Poincar{\'e} group extended by dilations. The causal structure provides the Poincar{\'e} sector, and the single unit (the clock) fixes the scale.

Finally, regarding the role of $c$: kinematically, MPSV~\cite{matsas-2024} correctly identify $c$ as a mere unit converter between the fundamental time and the derivative space. Yet, dynamically, $c$ retains a privileged status. It is the unique velocity parameter that ensures the form invariance of the equations of motion, such as those of electrodynamics~\cite{svozil-relrel,svozil-2001-convention}. Thus, while the clock fixes the scale of the universe, $c$ defines the rigid phenomenal structure necessary for the clock to operationally tick in such a way as to guarantee the form invariance of the physical laws under Lorentz boosts.


\begin{acknowledgments}
This text was partially created and revised with assistance from one or more of the following large language models:  Gemini 3 Pro
and gpt-5.1-high. All content, ideas, and prompts were provided by the authors.

This research was funded in whole or in part by the Austrian Science Fund (FWF) Grant \href{https://doi.org/10.55776/PIN5424624}{10.55776/PIN5424624}.
The author acknowledges TU Wien Bibliothek for financial support through its Open Access Funding Programme.

The authors declare no conflict of interest.
\end{acknowledgments}

\bibliography{svozil}


\end{document}
